\section{Cross-Cutting Analyses}
\label{sec:crosscutting}

\subsection{Convergence Parity}
\label{sec:crosscutting:convergence}

A prerequisite for attributing throughput differences to hardware rather than
numerical divergence is that both platforms converge to the same loss trajectory
under identical seeds.  We verify this via Gate~2 of the verification harness
(\S\ref{sec:methodology}): a single-GPU dense pretrain is run on each vendor
with \texttt{setup\_seed(42)} and bf16 AMP; the cross-entropy loss at step~100
must agree within $\pm 2\%$ relative.

\begin{table}[h]
\centering
\caption{Convergence parity (dense pretrain, bf16, 1 GPU, seed=42).
$\delta_{\text{rel}} = |L_{\text{H100}} - L_{\text{MI350}}| / \max(L_{\text{H100}}, L_{\text{MI350}})$.
Tolerance: $\leq 2\%$.}
\label{tab:convergence}
\small
\begin{tabular}{lrrrr}
\toprule
\textbf{Step} & \textbf{H100 loss} & \textbf{MI350 loss} & $\delta_{\text{rel}}$ (\%) & \textbf{Pass?} \\
\midrule
100  & \PH{} & \PH{} & \PH{} & \PH{} \\
500  & \PH{} & \PH{} & \PH{} & \PH{} \\
1000 & \PH{} & \PH{} & \PH{} & \PH{} \\
\midrule
Final (epoch 2) & \PH{} & \PH{} & \PH{} & \PH{} \\
\bottomrule
\end{tabular}
\end{table}

PyTorch-ROCm aliases the \texttt{torch.cuda.*} namespace and uses the same
cuBLAS-equivalent kernel dispatch path for GEMM workloads via hipBLASLt, so
convergence parity is expected in bf16.  Any deviation beyond 2\% would
indicate a fused-kernel discrepancy (e.g.\ in the SwiGLU or RMSNorm fused
backward) and would require per-operator investigation before the
throughput numbers can be meaningfully compared.

For the MoE configuration, we additionally verify that the router's soft-max
probabilities follow the same distribution at initialization (both vendors use
PyTorch's default Kaiming-uniform initialization, which must be seeded
identically).  The \texttt{aux\_loss} trajectory is a proxy for router load
balance and is reported alongside cross-entropy loss in Table~\ref{tab:convergence_moe}.

\begin{table}[h]
\centering
\caption{MoE convergence parity (bf16, 1 GPU, seed=42).  Both CE loss and
auxiliary load-balance loss are compared.}
\label{tab:convergence_moe}
\small
\begin{tabular}{lrrrr}
\toprule
\textbf{Step} & \textbf{H100 CE / aux} & \textbf{MI350 CE / aux} & $\delta_{\text{rel}}^{\text{CE}}$ (\%) & \textbf{Pass?} \\
\midrule
100  & \PH{} / \PH{} & \PH{} / \PH{} & \PH{} & \PH{} \\
500  & \PH{} / \PH{} & \PH{} / \PH{} & \PH{} & \PH{} \\
1000 & \PH{} / \PH{} & \PH{} / \PH{} & \PH{} & \PH{} \\
\bottomrule
\end{tabular}
\end{table}

\subsection{dtype Maturity: bf16, fp16, and the FP8 Outlook}
\label{sec:crosscutting:dtype}

\paragraph{bf16 (primary).}
Both H100 and MI350 natively support bf16 tensor cores.  MiniMind trainers
default to bf16 via \texttt{torch.amp.autocast(device\_type=device, dtype=torch.bfloat16)}.
The GradScaler is instantiated but effectively a no-op for bf16 (loss scaling
is not needed due to the wider dynamic range).  We observe numerically stable
training on both platforms with no manual loss-scale intervention.

\paragraph{fp16 (secondary).}
MiniMind supports fp16 via the same AMP path with GradScaler enabled.  fp16
is more sensitive to gradient underflow, particularly in the early warmup
phase of pretraining.  Table~\ref{tab:pretrain_dense} includes fp16 rows so
readers can verify that neither vendor exhibits systematic instability beyond
the known GradScaler overhead.

\paragraph{FP8 (future).}
Neither the current training scripts nor \texttt{torch.amp} expose a stable
FP8 path in PyTorch~2.10.  NVIDIA exposes FP8 via Transformer Engine (TE);
AMD exposes a parallel path via ROCm~7.x \texttt{hipBLASLt} FP8 GEMM
routines.  Because MiniMind's hand-written layers bypass TE's
\texttt{Linear} wrappers, enabling FP8 would require manual \texttt{Float8Linear}
substitution and is left as future work.  We flag in
Table~\ref{tab:dtype_matrix} which dtypes were exercised and which are
expected to become available.

\begin{table}[h]
\centering
\caption{dtype availability and training-stability status (MiniMind~2.10 baseline).}
\label{tab:dtype_matrix}
\small
\begin{tabular}{lllll}
\toprule
\textbf{dtype} & \textbf{H100 support} & \textbf{MI350 support} & \textbf{Trainer path} & \textbf{Status} \\
\midrule
bf16  & \checkmark & \checkmark & \texttt{autocast bf16 + no-op scaler} & Stable \\
fp16  & \checkmark & \checkmark & \texttt{autocast fp16 + GradScaler}   & Stable \\
fp32  & \checkmark & \checkmark & \texttt{nullcontext (default CPU)}    & Baseline only \\
FP8   & TE path    & ROCm 7.x   & Not implemented                      & Future \\
\bottomrule
\end{tabular}
\end{table}

\subsection{MoE Routing Behavior}
\label{sec:crosscutting:moe}

MiniMind's MoE FFN routes each token to $k{=}1$ of $N{=}4$ routed experts
plus a shared expert that always runs.  The router is a linear projection
followed by softmax; the top-$k$ gate values weight the expert outputs.
A load-balancing auxiliary loss $\mathcal{L}_{\text{aux}} = \alpha \sum_i f_i \cdot P_i$
(fraction of tokens $f_i$ vs.\ mean gate probability $P_i$ per expert) is
added to the primary CE loss at every step.

The key vendor-specific concern is the \emph{token-dispatch scatter/gather}
pattern: on CUDA, expert computation can be fused using custom Triton kernels
(or cuBLAS batched-GEMM with grouped indices); on ROCm, the same
\texttt{torch.index\_select} / \texttt{scatter\_add} path is used, falling
back to hipBLASLt for the GEMM portion.  For MiniMind's small expert count
($N{=}4$) and small hidden dimension (768), the dispatch overhead is
negligible relative to the GEMM time, but we track it explicitly via
step-time decomposition.

\begin{table}[h]
\centering
\caption{MoE expert utilization at epoch 1 end (pretrain, bf16, 1 GPU).
An ideal load-balanced router assigns $25\%$ of tokens to each routed expert.
Utilization $> 5\%$ deviation from $25\%$ may indicate router collapse.}
\label{tab:moe_balance}
\small
\begin{tabular}{lrrrrrr}
\toprule
\textbf{Vendor} & \textbf{Expert 0 (\%)} & \textbf{Expert 1 (\%)} &
\textbf{Expert 2 (\%)} & \textbf{Expert 3 (\%)} &
\textbf{Shared (\%)} & \textbf{aux\_loss} \\
\midrule
H100  & \PH{} & \PH{} & \PH{} & \PH{} & 100 & \PH{} \\
MI350 & \PH{} & \PH{} & \PH{} & \PH{} & 100 & \PH{} \\
\bottomrule
\end{tabular}
\end{table}

\subsection{\texttt{torch.compile} Codegen Quality}
\label{sec:crosscutting:compile}

\texttt{torch.compile} is exercised via the \texttt{--use\_compile 1} flag
available in all MiniMind trainers.  The compiler backend on CUDA (Ampere+)
is Triton/inductor, which generates efficient CUDA kernels for the attention
SDPA, layer norms, and element-wise operations.  On ROCm, the same
\texttt{torch.compile} path uses the Triton-ROCm backend (upstreamed into
PyTorch~2.3+), which emits LLVM-rooted HIP kernels via HSACO.

Known limitation: \texttt{torch.compile} is automatically disabled for LoRA
runs (\texttt{train\_lora.py:165}) because the monkey-patched
\texttt{Linear.forward} is not graph-capture-compatible.  This limitation
applies equally on both vendors.

\begin{table}[h]
\centering
\caption{\texttt{torch.compile} speedup over eager mode (pretrain, dense,
bf16, 8 GPUs, epoch 1 throughput).  Speedup is computed as
$\text{tok/s}_{\text{compile}} / \text{tok/s}_{\text{eager}} - 1$.}
\label{tab:compile}
\small
\begin{tabular}{lrrrr}
\toprule
\textbf{Vendor} & \textbf{Eager tok/s} & \textbf{Compile tok/s} & \textbf{Speedup (\%)} & \textbf{Compile time (s)} \\
\midrule
H100  & \PH{} & \PH{} & \PH{} & \PH{} \\
MI350 & \PH{} & \PH{} & \PH{} & \PH{} \\
\bottomrule
\end{tabular}
\end{table}

\PH{Discuss observed compile-time and warm-up overhead on each vendor once
data are available.  On AMD, the first-run Triton-ROCm kernel compilation
can take 5--15 minutes for a new model shape, which inflates the first-epoch
wall time; subsequent runs use the HSACO cache (\texttt{\$HOME/.triton/cache}).}

\subsection{NCCL vs.\ RCCL: DDP All-Reduce Bandwidth}
\label{sec:crosscutting:collective}

MiniMind uses PyTorch DDP, which performs one all-reduce per parameter tensor
per backward pass.  For a 64M-parameter dense model in bf16, the total
gradient buffer size is $64{\times}10^6 \times 2\,\text{B} \approx 128\,\text{MB}$
per all-reduce.  For the 198M-A64M MoE model (total parameters, bf16), the
buffer is $\approx 396\,\text{MB}$.

On the H100 side, NCCL uses NVLink~4 (900\,GB/s bidirectional per GPU pair
via NVSwitch) for intra-node all-reduce; effective ring-allreduce bandwidth
scales with the number of GPUs up to the bisection bandwidth limit.  On the
MI350 side, RCCL uses AMD Infinity Fabric (the specific bandwidth for MI350
in an 8-GPU OAM configuration is \PH{}) exposed as the \texttt{"nccl"} backend
(RCCL is API-compatible with NCCL and is aliased transparently by
\texttt{dist\_backend()} in \texttt{trainer/device\_utils.py}).

\begin{table}[h]
\centering
\caption{DDP all-reduce wall time for a 128\,MB gradient buffer (bf16, 8 GPUs,
bus-bandwidth measurement via \texttt{nccl-tests} or equivalent).}
\label{tab:allreduce}
\small
\begin{tabular}{lrrrr}
\toprule
\textbf{Vendor} & \textbf{Backend} & \textbf{Latency ($\mu$s)} &
\textbf{Bandwidth (GB/s)} & \textbf{Peak (\%)} \\
\midrule
H100  & NCCL  & \PH{} & \PH{} & \PH{} \\
MI350 & RCCL  & \PH{} & \PH{} & \PH{} \\
\bottomrule
\end{tabular}
\end{table}

Because MiniMind's model is small relative to the intra-node interconnect
bandwidth on either platform, we expect all-reduce to be a sub-dominant
cost ($< 5\%$ of step time at 8 GPUs) for both vendors.  This simplifies
the strong-scaling analysis in \S\ref{sec:training:pretrain}: near-linear
scaling is driven by compute efficiency, not communication overhead.
